
\usepackage[english]{babel}
\usepackage{iftex}
\ifPDFTeX
\usepackage[T1]{fontenc}
\usepackage[utf8]{inputenc}
\usepackage{lmodern}
\else
\usepackage{fontspec}
\defaultfontfeatures{Ligatures=TeX}
\setmainfont{Latin Modern Roman}
\setsansfont{Latin Modern Sans}
\setmonofont{Latin Modern Mono}
\usepackage{newunicodechar}
\newfontfamily\unicodeSymbolFont{Segoe UI Symbol}[
  BoldFont=Segoe UI Symbol,
  ItalicFont=Segoe UI Symbol,
  BoldItalicFont=Segoe UI Symbol
]
\newunicodechar{⟂}{{\unicodeSymbolFont ⟂}}
\fi
\usepackage{amsmath,amsfonts}
\usepackage{booktabs,multirow,xcolor,colortbl}
\usepackage{bm}%to use a more powerful version of \boldsymbol
\usepackage{fvextra}
\usepackage{csquotes}
\usepackage{textcomp}
\ifPDFTeX
\DeclareUnicodeCharacter{27C2}{\ensuremath{\perp}}
\fi

\newcommand{\AllSample}{ \mathcal Y^T }

\usepackage[citestyle=authoryear-comp,bibstyle=../Preamble/mystyle_lowercase,maxbibnames=5,minbibnames=1,maxnames=2, minnames=1,datezeros=false,date=long,isbn=false,url=false,backend=biber,useprefix=true]{biblatex}
\addglobalbib{../Preamble/Macro_database_biblatex.bib}	

\usepackage{upquote}
\usepackage{ragged2e}
\usepackage{microtype}

\usepackage{multimedia}
\usepackage{graphicx}
\usepackage[export]{adjustbox}
\usepackage{epstopdf}
\usepackage[labelfont=bf,format=hang, figurename=Fig.]{caption}
\usepackage[subrefformat=parens,labelformat=empty]{subcaption}
\usepackage{pgfpages} % for multiple-screen presentations

\usepackage{array}

\newcolumntype{L}[1]{>{\RaggedRight\let\newline\\\arraybackslash\hspace{0pt}}m{#1}}
\newcolumntype{C}[1]{>{\centering\let\newline\\\arraybackslash\hspace{0pt}}m{#1}}
\newcolumntype{R}[1]{>{\raggedleft\let\newline\\\arraybackslash\hspace{0pt}}m{#1}}
%https://tex.stackexchange.com/questions/12703/how-to-create-fixed-width-table-columns-with-text-raggedright-centered-raggedlef


%%% use serif font for math
\usefonttheme[onlymath]{serif}

%%% define blue highlighting
\newcommand{\highlight}[1]{{\color{blue}{#1}}}

%%%% define command for setting sources

\newcommand{\Source}{}

%%%
\beamertemplatenavigationsymbolsempty


%%% define themes
\usetheme{Boadilla}
\usecolortheme{rose}
\usecolortheme[named=blue]{structure}

\newcommand\textbox[1]{%
  \parbox{.333\textwidth}{#1}%
}

\setbeamertemplate{sections/subsections in toc}[sections numbered]


%%%% set coloring in  foot and head
\setbeamercolor{alerted text}{fg=red!80!black}
\setbeamertemplate{enumerate items}[default] % numbers enumerate items without encircling them
\setbeamertemplate{itemize items}[ball] % numbers enumerate items without encircling them
\setbeamercolor{section in head/foot}{fg=black, bg=white}
\setbeamercolor{subsection in head/foot}{fg=black, bg=white}
\setbeamercolor{title}{fg=black}
\setbeamercolor{subtitle}{fg=blue}

\setbeamertemplate{caption}[numbered]

%%% set overlays as default
\beamerdefaultoverlayspecification{<+->}

%%%% settings for output mode

\mode<beamer>
{
%\setbeameroption{show notes on second screen=right}
}

\mode<trans>
{\renewcommand{\highlight}[1]{{\textbf{#1}}}}


\mode<handout>
{
\renewcommand{\highlight}[1]{{\textbf{#1}}}
%\pgfpagesuselayout{4 on 1}[a4paper,border shrink = 5mm,landscape]
%\setbeamercolor{background canvas}{bg=black!5}
}


%%% define Exercise counter
\newcounter{Exercise}
\resetcounteronoverlays{Exercise}
\setcounter{Exercise}{0}
\newcommand{\Exercise}{\refstepcounter{Exercise}\emph{\structure{Übung \theExercise}: }}

%%%% define environments for theorems etc and make them numbered
\setbeamertemplate{theorems}[numbered]

\newtheorem{assum}{Assumption}
\setbeamertemplate{assum}[numbered]

\newtheorem{defin}{Definition}
\setbeamertemplate{defin}[numbered]

\newtheorem{propos}{Proposition}
\setbeamertemplate{propos}[numbered]


%%%% Make Parts show up in TOC
\makeatletter
\AtBeginPart{%
  \addtocontents{toc}{\protect\beamer@partintoc{\the\c@part}{\beamer@partnameshort}{\the\c@page}}%
}
%% number, shortname, page.
\providecommand\beamer@partintoc[3]{%
  \ifnum\c@tocdepth=-1\relax
    % requesting onlyparts.
    \makebox[6em]{#1.} #2
    \par
  \fi
}
\define@key{beamertoc}{onlyparts}[]{%
  \c@tocdepth=-1\relax
}
\makeatother%

%%%% fix bug in beamer package with note itemize
%%% http://tex.stackexchange.com/questions/1480/changing-the-textwidth-of-the-notes-in-beamer-repost-from-so
\makeatletter
\defbeamertemplate{note page}{infolines}
{%
  \vskip.75em
  \nointerlineskip
  \begin{minipage}{\textwidth} % this is an addition
  \footnotesize \insertnote
  \end{minipage}               % this is an addition
}
\makeatother

\setbeamertemplate{note page}[infolines]


\newcommand{\backupbegin}{
   \newcounter{framenumberappendix}
   \setcounter{framenumberappendix}{\value{framenumber}}
}
\newcommand{\backupend}{
   \addtocounter{framenumberappendix}{-\value{framenumber}}
   \addtocounter{framenumber}{\value{framenumberappendix}}
}

\hypersetup{bookmarksnumbered=true,naturalnames=true,pdfhighlight=/N,citebordercolor={1 1
1},linkbordercolor={1 1 1},colorlinks=true,anchorcolor=black,linkcolor=blue,citecolor=green,urlcolor=green,
pdfpagemode=UseOutlines,breaklinks=true,
pdfkeywords=Dynare, pdfsubject=Dynare,pdfstartpage=1} %pdfpagemode=fullpage

\usepackage[copyright]{ccicons}

\usepackage{tikz,pgfplots}
\pgfplotsset{compat=1.18}
\usetikzlibrary{arrows,calc,decorations.pathreplacing,decorations.shapes}
\pgfdeclarelayer{background}
\pgfdeclarelayer{foreground}
\pgfsetlayers{background,main,foreground}

\usepackage{relsize}
\newcommand\LM{\ensuremath{\mathit{LM}}}
\newcommand\IS{\ensuremath{\mathit{IS}}}
\DeclareMathOperator{\std}{std}


\setbeamertemplate{footline}{
\begin{beamercolorbox}{subsection in head/foot}
\insertsectionnavigationhorizontal{.95\textwidth}{\hskip0pt
\hfill}{\hfill \insertframenumber /\inserttotalframenumber}
\end{beamercolorbox}
}

\newenvironment{witemize}{\itemize\addtolength{\itemsep}{8pt}}{\enditemize}
\newenvironment{wenumerate}{\enumerate\addtolength{\itemsep}{8pt}}{\endenumerate}

\usepackage[]{minted}
\usepackage{tcolorbox}
\usepackage{etoolbox}
\BeforeBeginEnvironment{minted}{\begin{tcolorbox}[boxsep=0mm]}%
\AfterEndEnvironment{minted}{\end{tcolorbox}}%

\AtBeginEnvironment{minted}{%
  \renewcommand{\fcolorbox}[4][]{#4}}  